Nessa seção são apresentados os principais modelos matemáticos utilizados para representar o ambiente simulado e os aspectos relacionados ao descarregamento de tarefas, englobando modelos de comunicação celular, consumo energético e função objetivo. Para facilitar a leitura e entendimento é apresentada a Tabela \ref{tab:notations} com um resumo das principais notações utilizadas.

\begin{table}[ht]
\centering
\caption{Resumo das Notações Utilizadas}
\label{tab:notations}
{\scriptsize
\begin{tabularx}{\columnwidth}{l X}
\toprule
\textbf{Símbolo} & \textbf{Descrição} \\
\midrule
\multicolumn{2}{l}{\textit{Modelo de Rede e Sistema}} \\
$v_0$ & Veículo cliente (gerador de tarefas) \\
$\mathcal{V}(t)$ & Conjunto de veículos vizinhos no tempo $t$ \\
$B_{\mathrm{V2I}}, B_{\mathrm{V2V}}$ & Largura de banda de uplink (V2I) e sidelink (V2V) \\
$Q_{\mathrm{RSU}}, Q_v$ & Capacidade da fila no RSU e no veículo $v$ \\
\midrule
\multicolumn{2}{l}{\textit{Modelo de Tarefa}} \\
$\tau_j$ & Tarefa $j$ definida por $(C_j, D_j, s_j)$ \\
$C_j$ & Complexidade de CPU da tarefa $j$ (ciclos) \\
$D_j$ & Prazo (deadline) da tarefa $j$ (segundos) \\
$s_j$ & Tamanho dos dados de entrada da tarefa $j$ (bits) \\
$m$ & Modo de execução ($0$: local, $1$: RSU, $2$: V2V) \\
$f^{(m)}$ & Frequência de CPU no modo $m$ \\
\midrule
\multicolumn{2}{l}{\textit{Métricas de Desempenho}} \\
$\mathrm{JCT}_j^{(m)}$ & Tempo de conclusão (Job Completion Time) no modo $m$ \\
$t_j^{\mathrm{tx},(m)}$ & Tempo de transmissão (upload) dos dados \\
$t_j^{\mathrm{fila},(m)}$ & Tempo de espera na fila do destino \\
$E_j^{(m)}$ & Energia consumida pela tarefa $j$ no modo $m$ \\
$\alpha_c, V, P_{\mathrm{TX}}$ & Parâmetros físicos (capacitância, tensão, potência TX) \\
\midrule
\multicolumn{2}{l}{\textit{Objetivo de Otimização}} \\
$\Phi_j$ & Contribuição de QoS da tarefa $j$ (função de custo) \\
$\alpha, \beta, \gamma$ & Pesos de latência, energia e penalidade de prazo \\
$\hat{\tau}$ & Taxa de cumprimento de prazos (Success Rate) \\
\bottomrule
\end{tabularx}}
\end{table}

\subsection{Modelo de Rede}

Consideramos um ambiente \acrotip{VEC} heterogêneo em um cenário urbano, composto por uma \acrotip{RSU} fixa e um conjunto de veículos vizinhos $\mathcal{V}(t)$ que se deslocam segundo um modelo de mobilidade urbana. O veículo cliente $v_0$ dispõe de duas interfaces de comunicação para descarregamento:

\begin{itemize}[nosep,leftmargin=*]
    \item \textbf{V2I (Vehicle-to-Infrastructure):} Utiliza tecnologia \acrotip{5G} NR (banda n78) com frequência portadora de \SI{3{,}5}{\giga\hertz} e largura de banda $B_{\mathrm{V2I}} = \SI{20}{\mega\hertz}$. O modelo de perda de percurso segue a norma 3GPP TR 38.901 UMa LoS, considerando a altura da \acrotip{RSU} ($h_{\mathrm{BS}}=\SI{10}{\meter}$) e dos veículos ($h_{\mathrm{UT}}=\SI{1{,}5}{\meter}$).
    \item \textbf{V2V (Vehicle-to-Vehicle):} Utiliza sidelink PC5 (NR-V2X) operando em \SI{5{,}9}{\giga\hertz} com largura de banda $B_{\mathrm{V2V}} = \SI{20}{\mega\hertz}$. A propagação é modelada conforme a norma 3GPP TR 37.885 Urban LoS.
\end{itemize}

A taxa de transmissão instantânea para ambos os links é obtida pela capacidade de Shannon, $R = B \log_2(1 + \mathrm{SNR})$, onde a relação sinal-ruído ($\mathrm{SNR}$) integra a potência de transmissão ($P_{\mathrm{TX}} = \SI{0{,}2}{\watt}$), o ganho do canal e o ruído térmico AWGN calculado com base na temperatura de ruído (\SI{290}{\kelvin}) e na figura de ruído do receptor (\SI{5}{\decibel} para V2I e \SI{9}{\decibel} para V2V). A \acrotip{RSU} e cada veículo $v \in \mathcal{V}(t)$ gerenciam filas de serviço com capacidade $Q_{\max} = 25$ para o processamento de tarefas pendentes.

\subsection{Modelo de Tarefa}

Cada tarefa $\tau_j = (C_j, D_j, s_j)$ é caracterizada por sua complexidade
de CPU $C_j$ (ciclos), prazo $D_j$ (segundos) e tamanho de dados de entrada
$s_j$ (bits).
Três modos de execução estão disponíveis:
\begin{itemize}[nosep,leftmargin=*]
  \item \textbf{Local ($m=0$):} executado em $v_0$ com frequência $f_{\mathrm{loc}}$.
  \item \textbf{\acrotip{RSU} ($m=1$):} carregado ao \acrotip{RSU} e executado em $f_{\mathrm{RSU}}$.
  \item \textbf{\acrotip{V2V} ($m=2$):} descarregado a um veículo vizinho $v^* \in \mathcal{V}(t)$.
\end{itemize}
Um modo $m$ de descarregamento é viável se o tempo de contato com o destino for suficiente para a transmissão dos dados de entrada e recepção dos resultados.

\subsection{Modelo de Latência}

O Tempo de Conclusão de Tarefa do modo $m$ para a tarefa $\tau_j$ é
\begin{equation}
  \mathrm{JCT}_j^{(m)} = t_j^{\mathrm{tx},(m)} + t_j^{\mathrm{fila},(m)}
                         + \frac{C_j}{f^{(m)}},
  \label{eq:jct}
\end{equation}
onde $t_j^{\mathrm{tx},(m)}$ é o tempo de transmissão (upload) e
$t_j^{\mathrm{fila},(m)}$ é o tempo de espera residual na fila do destino,
estimado como o tempo de execução restante das tarefas já na fila.

\subsection{Modelo de Energia}

A energia consumida pelo modo $m$ é
\begin{equation}
  E_j^{(m)} = \alpha_c^{(m)} \bigl[V^{(m)}\bigr]^2 C_j
             + P_{\mathrm{TX}}^{(m)}\,t_j^{\mathrm{tx},(m)},
  \label{eq:energy}
\end{equation}
onde $\alpha_c^{(m)}$ é a capacitância do chip, $V^{(m)}$ é a tensão de
alimentação na frequência $f^{(m)}$ e $P_{\mathrm{TX}}^{(m)}$ é a potência
de transmissão.
Para o modo local, o termo de transmissão é nulo.

\subsection{Objetivo de Otimização}

A contribuição de \acrotip{QoS} por tarefa é
\begin{equation}
  \Phi_j = \alpha\,\mathrm{JCT}_j + \beta\,E_j
            + \gamma\,\mathbf{1}[\mathrm{JCT}_j > D_j],
  \label{eq:obj}
\end{equation}
com parâmetros configuráveis de peso (por padrão $\alpha=\beta=0{,}35$) e uma constante de penalidade por violação de prazo $\gamma=0{,}015$.
O objetivo do sistema é minimizar $\sum_j \Phi_j$ maximizando a
taxa de cumprimento de prazo
$\hat{\tau} = |\{j : \mathrm{JCT}_j \le D_j\}|/N$.
